\textbf{Время:} резерв

\textbf{Тезис:} инвариантный относительный импульс отличается от пространственной разности в произвольной системе.

\textbf{Визуал:} полная разность и полусумма; система покоя пары рядом с LCMS.

\textbf{Текст:} Четырёхимпульс содержит энергию и обычный импульс. Здесь q — полная разность, K — полусумма, метрика имеет знаки плюс, минус, минус, минус. В обзоре Lisa относительный импульс для одинаковых масс вдвое меньше нашего.

Q-inv определяется как корень из пространственной разности в квадрате минус разность энергий в квадрате. Для равных масс в системе покоя пары энергии одинаковы, и q-inv равен удвоенному модулю импульса одной частицы. В LCMS поперечное движение пары остаётся, поэтому системы не совпадают.

Средний поперечный импульс определяется через векторную полусумму. Изменение определения q без пересчёта фазы и экспоненты меняет масштаб извлечённого радиуса.

\textbf{Источники и условия:} Метрика \((+,-,-,-)\). У Lisa для равных масс \(q_{\rm Lisa}=q/2\). Определения \(q\) и \(k_T\) совпадает со статьёй Краевой (2024).
